% Pursuing Stacks (cotes 134-2 to 134-8): archive cote 134-6, site page 120 (PDF page 121), typescript page 438.
% typescript. Draft compared: ps5.tex:7788-7872.
% AI-transcribed by Claude Opus 5.5 (claude-opus-5-5, adaptive thinking, effort high), Message Batches API, batch ps3 (msgbatch_01XdFmuVVBAcgc7EFxdtCqsj), 2026-11-12; not yet checked by eye.
\scanpage{134-6/120}
\mspage{438}

which may be computed using injective resolutions of the second argument
$F^\bullet$ in $R\underline{Hom}_k(L.,\ F^\bullet)$ (but not, in general, by using projective resolu-%
tions of L.), i.e.
\[ (133) \qquad R\underline{Hom}_k(L.,F^\bullet) \hand{$\simeq$} \underline{Hom}_k^{\bullet\bullet}(L.\ ,\hand{$\mathcal{F}$}^\bullet) \ , \]
where $\hand{$\mathcal{F}$}^\bullet$ is an injective resolution of $F^\bullet$ (i.e. a complex in $A^\wedge_{\hand{$\mathcal{M}$}}$ with
degrees bounded from below and \over{\unsure{pro}}{in}jective components, endowed with a
quasi-isomorphism $F.\to\hand{$\mathcal{F}$}.$ ).

Applying now $R\Gamma_A$ to both members of (133), and using the similar
isomorphism for $RHom_k$ (valid by cor.2), we get the familiar formula
\[ (134) \qquad RHom_k(L.,F^\bullet) \hand{$\simeq$} RH^\bullet(R\underline{Hom}_k(L.,F^\bullet)) \ , \]
(where in accordance with (124), we wrote $RH^\bullet$ instead of $R\Gamma_A$ ), where
however $F^\bullet$ is \inserted{now} a complex of \xout{xxaxxx} pres\xout{x}heaves with values in a
k-additive category \hand{$\mathcal{M}$} (satisfying the conditions dual to a) to d) in
cor.1)\hand{, not just a presheaf of k-modules.}

\margin{\unsure{$\vdash$}}Replacing \hand{$\mathcal{M}$} by a category \hand{$\mathcal{N}$} satisfying the assumptions \deleted{of} \inserted{dual to those of} cor.1 ,
we get likewise a total \inserted[typed]{left} derived functor
\[ (135) \qquad (F.,L.')\mapsto F \overset{L}{\over{\ast}{\otimes}}_k L.' \ : \quad D^-(A^\wedge_{\hand{$\mathcal{N}$}}) \text{ x } D^-(B^\wedge_k) \qquad D^-(A^\wedge_{\hand{$\mathcal{N}$}}) \ , \]
which can be defined  using projective resolutions of the first argument
F. (but not by using projective resolutions of L.' ), giving rise to
the isomorphism (dual to (134))
\[ (134') \qquad F.\ \overset{L}{\ast}_k L.' \simeq LH.(F.\ \overset{L}{\over{\ast}{\otimes}}_k L.') \qquad . \]
The duality relationship between $R\underline{Hom}_k$ and $\overset{L}{\over{\ast}{\otimes}}_k$ can be expressed by
two obvious formulae, similar to (120) and (120') for $RHom_k$ and $\overset{L}{\ast}_k$ ,
which we leave to the reader.

\uline{Remarks} \hand{1)} If we want to consider ``multiplicative structure'' in the purely
``k-linear'' set-up, where the data is a small k-additive category P,
rather than a small category A (giving rise to $P = Add_k(A)$), and corres
-ponding cohomology \xout{structure} and homology operations $R\underline{Hom}_k$ and $\overset{L}{\over{\ast}{\otimes}}_k$ ,
not only $RHom_k$ and $\overset{L}{\ast}_k$ , the natural thing to do, it seems, is to
introduce a\hand{/}``\uline{diagonal map}''
\[ \hand{(*)} \qquad P \to P\otimes_k P \qquad , \]
where the tensor product in the second member \xout{is} can be defined in a
rather evident way (as solution of the obvious 2-universal problem in
terms of k-biadditive functors on PxP ), which will give rise in the
``usual'' way to a tensor product operation in both $P^{\hand{$\&$}}$ \hand{and $Q^{\&}$ (where}

\margin{120}
